\section*{第 \ref{ch:g12:projectile-motion} 章——自由落体与抛体运动}

\begin{solution}{exo:g12:projectile-motion:1}
$h = \tfrac12 \times 9.81 \times 2.0^2 = \qty{19.6}{m}$；
$v = gt = \qty{19.6}{m/s} \approx \qty{71}{km/h}$。
\end{solution}

\begin{solution}{exo:g12:projectile-motion:2}
月球上没有空气：二者都只受\omterm{def:g10:universal-gravitation:weight}{重力}，$m\vect a = m\vect g$ 给出
对任意质量都有 $\vect a = \vect g$。在地球上，羽毛所受的阻力
与其微小的\omterm{def:g10:universal-gravitation:weight}{重力}相当：它以约 \qty{1}{m/s} 的终端速率下落，
而铁锤几乎不受空气影响。
\end{solution}

\begin{solution}{exo:g12:projectile-motion:3}
$t = \sqrt{2h/g} = 0.45$、$1.4$、\qty{4.5}{s}；
$v = gt = 4.4$、$14$、\qty{44}{m/s}。时间随 $\sqrt h$ 增长：
高度增至 $\times 100$ 时，时间只增至 $\times 10$。
\end{solution}

\begin{solution}{exo:g12:projectile-motion:4}
$h = v_0^2/2g = \qty{11.5}{m}$，出现在 $t = v_0/g = \qty{1.53}{s}$ 时；
在 \qty{3.06}{s} 时回到原处，速率再次为 \qty{15}{m/s}（对称性）。
\end{solution}

\begin{solution}{exo:g12:projectile-motion:5}
$t = \sqrt{2 \times 0.80/9.81} = \qty{0.40}{s}$；
$x = 2.5 \times 0.40 = \qty{1.0}{m}$。不对：
下落时间只由竖直运动决定（独立性）。
\end{solution}

\begin{solution}{exo:g12:projectile-motion:6}
$(\ang{15},\ang{75})$ 与 $(\ang{30},\ang{60})$ \omterm{def:g12:projectile-motion:range}{射程}相同；
\ang{45} \omterm{def:g12:projectile-motion:range}{射程}最远。互余的一对中，角度较大者\omterm{def:g12:projectile-motion:range}{飞行时间}更长：
$t_f = 2v_0\sin\alpha/g$ 随 $\alpha$ 增大而增大。
\end{solution}

\begin{solution}{exo:g12:projectile-motion:7}
由 $4.905\,t^2 + 5.0\,t - 20 = 0$ 可得
$t = (-5 + \sqrt{25 + 392.4})/9.81 = \qty{1.57}{s}$；
$v = 5.0 + 9.81 \times 1.57 = \qty{20.4}{m/s}$。用\omterm{def:g10:energy-conservation:energy}{能量}法：
$v = \sqrt{5.0^2 + 2 \times 9.81 \times 20} = \qty{20.4}{m/s}$ ——
结果相同。
\end{solution}

\begin{solution}{exo:g12:projectile-motion:8}
$t_f = 2 \times 40 \times \sin\ang{30}/9.81 = \qty{4.1}{s}$；
$h = (40\sin\ang{30})^2/(2 \times 9.81) = \qty{20.4}{m}$；
$R = 40^2 \sin\ang{60}/9.81 = \qty{1.4e2}{m}$。
\end{solution}

\begin{solution}{exo:g12:projectile-motion:9}
$R = 10.5^2\sin\ang{40}/9.81 \approx \qty{7.2}{m}$。未计入的部分：
\omterm{def:g11:forces-and-motion:com}{质心}起跳较高（人体有一定大小），落地时则较低、更靠前（双腿前伸），
而且人并非质点——
缺少的 \qty{1.7}{m} 来自这一几何构型，而不是抛物线本身。
\end{solution}

\begin{solution}{exo:g12:projectile-motion:10}
两球同时落地：二者 $v_y(0) = 0$，故下落时间相同，
$t = \sqrt{2 \times 1.2/9.81} = \qty{0.49}{s}$；被水平射出的球落在
距离 $300 \times 0.49 \approx \qty{1.5e2}{m}$ 处。该实验检验
水平运动与竖直运动的独立性。
\end{solution}

\begin{solution}{exo:g12:projectile-motion:11}
$\sin\bigl(2(\ang{90}-\alpha)\bigr) = \sin(\ang{180}-2\alpha) =
\sin(2\alpha)$：$R$ 相同。此处 $\sin(2\alpha) = gR/v_0^2 =
9.81 \times 50/625 = 0.785$：$\alpha = \ang{25.9}$ 或 \ang{64.1}。
\end{solution}

\begin{solution}{exo:g12:projectile-motion:12}
$R_{\max} = v_0^2/g = \qty{6.5}{m}$（取 \ang{45}）；
$h_{\max} = v_0^2/2g = \qty{3.3}{m}$（竖直向上）。
$\sin(2\alpha) = 9.81 \times 5.0/64 = 0.766$：$\alpha = \ang{25.0}$
或 \ang{65.0}。
\end{solution}

\begin{solution}{exo:g12:projectile-motion:13}
$y(15) = 15\tan\ang{35} - \dfrac{9.81 \times 15^2}
{2 \times 18^2 \cos^2\ang{35}} = 10.5 - 5.1 = \qty{5.4}{m}$：
小球越过高 \qty{4.0}{m} 的篱笆，约高出 \qty{1.4}{m}。
\end{solution}

\begin{solution}{exo:g12:projectile-motion:14}
(a) 阻力与\omterm{def:g10:universal-gravitation:weight}{重力}平衡：$80 \times 9.81 \approx \qty{7.8e2}{N}$。
(b) $v(t)$ 的斜率越来越小，最终趋于
\qty{55}{m/s}。(c) 真空中：
$\sqrt{2 \times 9.81 \times 2000} \approx \qty{2.0e2}{m/s}$，
是实际 \qty{9}{m/s} 的二十倍——
雨滴下落的几乎全程都以终端速率运动。
\end{solution}

\begin{solution}{exo:g12:projectile-motion:15}
在 $t^* = D/(v_0\cos\alpha)$ 时，飞镖的高度为
$v_0\sin\alpha\, t^* - \tfrac12 g t^{*2} = D\tan\alpha - \tfrac12 g
t^{*2} = H - \tfrac12 g t^{*2}$，恰好等于下落猴子的
$y = H - \tfrac12 g t^{*2}$。二者都比无\omterm{def:g10:universal-gravitation:weight}{重力}时的位置低 $\tfrac12 gt^2$，
对任意 $v_0$ 都会相遇。
\end{solution}

\begin{solution}{pb:g12:projectile-motion:1}
\textbf{1.} 只受\omterm{def:g10:universal-gravitation:weight}{重力}（忽略空气）；由\omterm{def:g10:inertia:force}{牛顿}第二定律：
$m\vect a = m\vect g$，故 $\vect a = (0,\,-g)$。

\textbf{2.} $v_x = v_0\cos\alpha$，$v_y = v_0\sin\alpha - gt$。

\textbf{3.} $x = v_0\cos\alpha\, t$，
$y = v_0\sin\alpha\, t - \tfrac12 gt^2$。

\textbf{4.} $t = x/(v_0\cos\alpha)$:
$y = x\tan\alpha - g x^2/(2v_0^2\cos^2\alpha)$.

\textbf{5.} $v_0^2 = \dfrac{g d^2}{2\cos^2\alpha\,(d\tan\alpha -
1.05)} = \dfrac{173.0}{0.826 \times 3.96} = 52.9$:
$v_0 = \qty{7.28}{m/s} \approx \qty{26}{km/h}$.

\textbf{6.} $t = d/(v_0\cos\alpha) = 4.20/4.68 = \qty{0.90}{s}$.

\textbf{7.} 最高点出现在 $t = v_0\sin\alpha/g = \qty{0.57}{s}$ 时，高度
$2.00 + (v_0\sin\alpha)^2/2g = 2.00 + 1.58 = \qty{3.58}{m}$;
$0.57 < 0.90$：球正在篮圈处下降。

\textbf{8.} $y(0.90) = 0.90\tan\ang{50} - 0.182 = \qty{0.89}{m}$
高于释放点，即距地面 \qty{2.89}{m}：
比指尖高 \qty{0.29}{m}。

\textbf{9.} $v_x = \qty{4.68}{m/s}$;
$v_y = 5.57 - 9.81 \times 0.90 = \qty{-3.24}{m/s}$;
$v = \qty{5.69}{m/s}$，速度方向斜向下，
与水平方向夹角 $\arctan(3.24/4.68) = \ang{34.7}$。

\textbf{10.} 篮圈是水平圆环：球必须从上方穿过其
所在平面。入射越陡，
沿速度方向看去，篮圈开口越大——容错余地越大。

\textbf{11.} $1.05/26.0 = 4\%$：可忽略不计。由 $y(t_f) = 0$ 得
$t_f\,(v_0\sin\alpha - \tfrac12 g t_f) = 0$，故
$t_f = 2v_0\sin\alpha/g$.

\textbf{12.} $R = v_0\cos\alpha\, t_f = v_0^2\sin(2\alpha)/g$;
在 $\alpha = \ang{45}$ 时最大。

\textbf{13.} $v_0 = \sqrt{gR} = \sqrt{9.81 \times 26.0} =
\qty{16.0}{m/s} \approx \qty{57}{km/h}$.

\textbf{14.} $R_{\max} = 17.5^2/9.81 = \qty{31.2}{m} >
\qty{26.0}{m}$：可行，尚有 \qty{5}{m} 的\omterm{def:g12:projectile-motion:range}{射程}余量。

\textbf{15.} $\sin(2\alpha) = 255.1/306.3 = 0.833$:
$\alpha = \ang{28.2}$ 或 \ang{61.8}；\omterm{def:g12:projectile-motion:range}{飞行时间}
为 $2v_0\sin\alpha/g = \qty{1.69}{s}$ 和 \qty{3.14}{s}，仅
平射的投篮能赶在 \qty{2.0}{s} 蜂鸣器响起之前完成。

\textbf{16.} $v_0^2 = \dfrac{9.81 \times 36}{2\cos^2\ang{75}\,
(6.0\tan\ang{75} - 1.05)} = \dfrac{353.2}{2.86} = 123.5$:
$v_0 = \qty{11.1}{m/s}$.

\textbf{17.} $y(5.55) = 5.55\tan\ang{75} - 18.26 = \qty{2.45}{m}$
高于释放点，即距地面 \qty{4.45}{m}：
比篮板上沿高 \qty{0.50}{m}。

\textbf{18.} 最高点 $2.00 + (11.1\sin\ang{75})^2/2g \approx
\qty{7.9}{m}$；\omterm{def:g12:projectile-motion:range}{飞行时间} $6.0/(11.1\cos\ang{75}) = \qty{2.1}{s}$——
整整两秒，球在 \qty{8}{m} 的高处滞空。

\textbf{19.} 阻力缩短每次投篮的\omterm{def:g12:projectile-motion:range}{射程}，并使下落段变得更陡；
效应随速度增大，故 \qty{17.5}{m/s} 的全场远投
受影响最大——其实际\omterm{def:g12:projectile-motion:range}{射程}比
\qty{31}{m} 短几米，侵蚀了容错余量。

\textbf{20.} \emph{压哨球的结论}：\qty{26.0}{m}
的投篮需要 \qty{16.0}{m/s} 的速度，手臂约能提供
\qty{17.5}{m/s}，\omterm{def:g12:projectile-motion:range}{射程}余量为 \qty{5}{m}（$\approx 20\%$），
空气阻力会削减这一余量，但不能将其抹去。可能——这正是
为何每个赛季总会有几次球投进。
\end{solution}
